\subsection{Signed exclusions at \texorpdfstring{$11\times5$}{11x5} and \texorpdfstring{$8\times7$}{8x7}}\label{sec:signed-exclusions}

\begin{theorem}\label{thm:signed-exclusions}
One has
\[
 \zr(11,5)=\zs(11,5)=37,\qquad \zr(8,7)=\zs(8,7)=38 .
\]
\end{theorem}

\begin{lemma}\label{lem:bases-11-5-8-7}
One has $z(11,5)=21$, and the maximum base is unique up to row and
column permutations: ten rows carry the ten column pairs once each, and
one row is a singleton.
One has $z(8,7)=22$, and there are exactly two classes of maximum bases:
\begin{itemize}
\item type~A replaces one triple of a Fano triple system by two of its
      three pairs;
\item type~B adds one singleton row to a Fano triple system.
\end{itemize}
\end{lemma}

\begin{proof}
With row degrees $d_r$, $C_4$-freeness gives
$\sum_r\binom{d_r}2\le\binom n2$.
\begin{itemize}
\item \emph{$11\times5$.} Since $11\ge\binom52$, Lemma~\ref{lem:cap} gives
  $|F|\le21$ and forces ten rows of degree two, one on each column pair,
  and one row of degree one. Column permutations are transitive on the
  singleton's column.
\item \emph{$8\times7$.} The inequality $2d\le3+\binom d2$ gives
  $2|F|\le24+21$, so $|F|\le22$. At $|F|=22$, $\sum_r(d_r-2)(d_r-3)/2\le1$,
  so every row has degree two or three, with at most one row of degree one
  or four. This leaves the degree profiles $2^2\,3^6$, $1\,3^7$ and
  $2^3\,3^4\,4$.
\end{itemize}
In the last profile, each triple meets the degree-four row in at most one
column, since the base has no $C_4$. So each of the four degree-three
rows contains at least two of the three columns outside the degree-four
row, and hence a pair of them. Two triples share at most one column, so
these four pairs are distinct. Only three such pairs exist, so this
profile is impossible.

For $1\,3^7$, the seven triples cover all $21$ column pairs, so they form
a Fano system. The remaining row is a singleton, which gives type~B.

For $2^2\,3^6$, the three column pairs missed by the six triples have even
degree at every point, so they form a triangle. Adding that triangle
completes a Fano system, and the two pair rows are two of its edges,
which gives type~A.

The Fano system is unique up to relabelling: its three triples through
one point may be labelled $012,034,056$; choosing the triple through
$1,3$ as $135$ forces the remaining triples $146,236,245$.
Its automorphism group is transitive on points and on incident
point--triple flags. Hence each type is a single class.
\end{proof}

\begin{proof}[Proof of Theorem~\ref{thm:signed-exclusions}]
\emph{Lower bounds.} Explicit recursive-line displays of lengths $37$ and
$38$ on maximum bases were checked by the reference closure and an
independent C++ closure in the recorded audit. The portable replay in
packet X115/X87 rechecks them with the reference closure.

\emph{Reduction to hole-free displays.} We have
\[
U(11,5)=\lfloor(55+21)/2\rfloor=38,\qquad U(8,7)=\lfloor(56+22)/2\rfloor=39 .
\]
In both cases $mn+z(m,n)$ is even. So by the cell count
$mn=|F|+2|D|+H$, a display of length $U$ has no hole. It therefore
suffices to show that none of the three maximum bases of
Lemma~\ref{lem:bases-11-5-8-7} carries a hole-free display satisfying
$\SG$.

\emph{Heredity.} Suppose $(F,D)$ satisfies $\SG$ and $S\subseteq D$. Then
the display $(F,S)$, in which the cells of $D\setminus S$ become holes,
also satisfies $\SG$ (Lemma~\ref{lem:heredity}). We use this in two
ways:
\begin{itemize}
\item every pair of a signed display is a \emph{candidate}, that is, a
      pair whose one-pair display is accepted;
\item a set $S$ of candidates whose display $(F,S)$ fails $\SG$ cannot lie
      in a signed display. Such an $S$ is a \emph{cut}.
\end{itemize}

\emph{Region domains.} Fix a set $R$ of rows:
\begin{itemize}
\item at $11\times5$, the four pair rows through the singleton's column;
\item on type~A, its two pair rows;
\item on type~B, the three triples through the singleton's column.
\end{itemize}
The \emph{pattern} of a display is its set of pairs with both cells in
$R$. By heredity it is itself an accepted display on candidates. An
exhaustive enumeration of all partial matchings of the free cells of $R$
lists every such pattern. Every accepted pattern ($4{,}968$, $1{,}449$
and $353$ of them) is mapped onto one of $229$, $190$ and $18$ listed
representatives by a base automorphism preserving $R$. Each map is checked
to preserve the base.

\emph{Propositional encoding.} For a base, take one variable $x_P$ for
each candidate pair $P$, and the following clauses:
\begin{itemize}
\item exactly one selected pair meets each free cell;
\item for every cut $S$ found, the clause $\bigvee_{P\in S}\neg x_P$.
\end{itemize}
Call this formula $B$. For each representative $d$, let $L_d$ be the unit
literals fixing the region pattern to $d$. A hole-free signed display,
moved by an automorphism so that its pattern is a representative,
satisfies $B\wedge L_d$ for that $d$. The formulas contain no
symmetry-breaking clauses; the identical-row swap of
Proposition~\ref{prop:four-signed} has no instance on these bases.
Hence it suffices to refute $B\wedge L_d$ for every representative $d$.
The final formula adds \emph{selector} clauses: a new variable $s_d$
for each $d$, the clauses $s_d\Rightarrow L_d$, and $\bigvee_d s_d$.

The cuts were generated by counterexample-guided search. Their
provenance is irrelevant to the proof, because every cut is re-checked.

\emph{Refutations.} Each refuted formula is a subset of $B\wedge L_d$: the unit clauses
$L_d$, the propagating clauses, and the clauses of $B$ not satisfied
by unit propagation from $L_d$. The checker verifies this subset relation line by
line.

For a few domains the refutation is split on the partner of one free cell
$c$ outside $R$. Since $B$ contains the clause $\bigvee_{P\ni c}x_P$,
refuting $B\wedge L_d\wedge x_P$ for every candidate $P\ni c$ refutes
$B\wedge L_d$.

Every refutation is a CaDiCaL LRAT proof. It is accepted by cake\_lpr, a
formally verified checker \cite{CakeLPR}, and separately by the LRAT checker of the
drat-trim distribution.

\emph{Audit.} Every clause of each final formula is classified:
\begin{itemize}
\item exactly-one clauses are rebuilt from the candidate list, and
      selectors from the representative list;
\item every other clause must be a negative clause on candidates, and its
      pair set is re-checked as a failed display, by the C++ closure and by
      the reference closure;
\item every excluded non-candidate pair is re-checked;
\item the region literals are rebuilt from the representatives.
\end{itemize}
The coverage enumeration and the candidate/non-candidate split were
repeated with the reference closure.
In addition, a random sample of $34$ domains and split cases, at least
eight for each base, was refuted again by kissat, and drat-trim accepted
all $34$ DRAT proofs.

The counts are as follows.
\begin{center}
\begin{tabular}{lrrrrl}
\toprule
Base&Candidates&Domains&Cut clauses&Split domains&Split cell cases\\
\midrule
$11\times5$&387&229&5{,}168{,}225&4&$4\times21$\\
$8\times7$ type A&326&190&946{,}334&1&16\\
$8\times7$ type B&312&18&589{,}720&1&15\\
\bottomrule
\end{tabular}
\end{center}
The cut counts refer to the final formula of each base. At $11\times5$ and
on type~A, $213$ and $189$ domains respectively were refuted from an
earlier formula of the same kind, which was audited in the same way
($3{,}956{,}309$ and $946{,}270$ cut clauses).

No hole-free signed display of length $U(m,n)$ therefore exists, so
$\zs(11,5)\le37$ and $\zs(8,7)\le38$. The hierarchy completes the proof.
\end{proof}

\begin{remark}
These are upper bounds for the signed criterion, not for the actual
square length. The actual values remain in the intervals
$37\le\zt(11,5)\le38$ and $38\le\zt(8,7)\le39$.

Together with the other finite values and the width-five formula of
Xi and Chang \cite{XiChangFive},
Theorem~\ref{thm:signed-exclusions} completes the width-five comparison
(Corollary~\ref{cor:five-and-six}). At $8\times7$
the common value $38$ lies one below the cell cap $U(8,7)=39$.

At $14\times4$ the remaining interval
$37\le\zr(14,4)\le\zs(14,4)\le38$ is the subject of
Problem~5.1 in \cite{CC}.
\end{remark}

\begin{proposition}\label{prop:nine-five}
One has $\zr(9,5)=\zs(9,5)=\zt(9,5)=31$.
\end{proposition}
\begin{proof}
Lemma~\ref{lem:cap} gives $z(9,5)\le18$. The following display has
18 singleton cells and 13 pairs, and its singleton set is $C_4$-free.
Equal numbers form one selected pair; bullets are singletons and the
circle is a hole. Packet BOUNDARY checks its recursive-line closure.
\[
\begin{array}{c|ccccc}
 &0&1&2&3&4\\\hline
0 & \bullet & \bullet & 1 & 4 & 5\\
1 & \bullet & 3 & \bullet & 7 & 4\\
2 & \bullet & 1 & 2 & \bullet & 9\\
3 & \bullet & 6 & 8 & 12 & \bullet\\
4 & 9 & \bullet & \bullet & 12 & 10\\
5 & 11 & \bullet & 13 & \bullet & 2\\
6 & 13 & \bullet & 11 & 8 & \bullet\\
7 & 10 & 5 & \bullet & \bullet & \circ\\
8 & 6 & 7 & \bullet & 3 & \bullet\\
\end{array}
\]
The base is maximum and $U(9,5)=\lfloor(45+18)/2\rfloor=31$, so the
hierarchy proves the assertion.
\end{proof}

\begin{corollary}\label{cor:five-and-six}
For every $m\ge1$, $\zr(m,5)=\zs(m,5)$. Consequently
$\zr(m,n)=\zs(m,n)$ for every rectangle with $\min(m,n)\le6$, except
possibly $14\times4$ and $4\times14$.
\end{corollary}

\begin{proof}
Both criteria are invariant under transposition. At width five, rows one
through four follow from the one-, two- and three-column values and the
$5\times4$ equality in \cite{LQ,XY}. At five rows the common value is $17$
\cite{LQ,XY}; at six rows it is $20$ by Theorem~\ref{thm:small} and
transposition. Theorem~\ref{thm:further-finite} gives $24,27,34$ at rows
seven, eight and ten. Proposition~\ref{prop:nine-five} gives $31$ at nine
rows. Eleven rows are covered by Theorem~\ref{thm:signed-exclusions},
and every $m\ge12$ by Theorem~\ref{thm:otherwidths}.
Widths one to three are in \cite{LQ}, width four is
Corollary~\ref{cor:four-exception}, and width six is Theorem~\ref{thm:six}.
This covers every rectangle with a side at most six, with the stated
exception.
\end{proof}

